\question{Does Flower Labs, the company, receive anything?}
\label{sec:flowerlabs}
\statusMethod

\begin{shortanswer}
No. In the intended deployment the SuperLink is self-hosted on consortium infrastructure and Flower
is open-source software running on that host; the company has no access path and receives nothing.
One caveat was worth closing rather than explaining: the Flower CLI writes a default entry pointing
at a Flower-operated endpoint into a freshly generated configuration file. It is inert unless
explicitly invoked, but it is one word on a command line, so its removal is now a documented step on
the practice-side go-live checklist.
\end{shortanswer}

\subsection{The configuration default, and why it was closed}

A newly generated \texttt{\$HOME/.flwr/config.toml} contains, among the connection profiles:

\begin{quote}\ttfamily
[superlink.supergrid]\\
address = "supergrid.flower.ai"
\end{quote}

Nothing routes there unless someone runs the deployment command naming that profile explicitly. But
the gap between a self-hosted federation and transmission of model updates to a third party should
not be a single word on a command line that a tired operator could type. The deployment procedure now
requires removing the entry on every operator and practice machine and pinning the default profile to
the consortium SuperLink, so that an omitted argument also cannot reach outside the consortium. It is
cheap to do and easy to forget, which is why it is on a checklist rather than in prose.

\subsection{The Letter of Intent}

The Projektantrag \cite{projektantrag} records a Letter of Intent from FlowerAI. It concerns
technical support and expertise. \textbf{It is not a data-processing arrangement and should not be
relied on as one.} If any future configuration were to route data through infrastructure Flower Labs
operates, that would be a new processing relationship requiring its own legal basis, and it is not
what is planned.
